% Note: БСП (блуждание без самопересечений) = SAW (self-avoiding walk).
% Below, we use SAW at first mention and thereafter.

\paragraph*{Problem statement.}

Critical phenomena in spin models on regular lattices have been studied in considerable numerical detail. 
 When moving to a dynamic geometry, the system’s behavior changes: the spin-bearing substrate itself fluctuates, which alters the mechanism of ordering. 
 A polymer substrate carrying magnetic spins can be modeled as a self-avoiding walk (SAW), where the polymer may undergo a structural reconfiguration. This structural (coil--globule) transition may compete with or synchronize with the magnetic transition. Our goal is to consolidate the currently fragmented results into a coherent picture, with particular emphasis on determining the order of the phase transitions.


\paragraph*{Relevance.}

In recent years, spin models on flexible substrates, especially on polymers, have attracted significant attention as a rich class of systems where geometry and magnetic ordering mutually influence each other. 
Studying spin models on dynamic structures clarifies how structural states affect magnetic phase transitions and whether qualitatively new system-level properties emerge that are absent on regular lattices.


\paragraph*{Problem Development.}

Spin models on polymer chains constitute a class where magnetic ordering and geometric fluctuations are coupled. On regular lattices, the universal properties of critical phenomena are well established \cite{cardy1996scaling,stanley1971introduction,baxter2013exactly,PELISSETTO2002549,kosterlitz1973ordering,pokrov}, yet transitioning to a dynamic geometry reshapes the phase diagram: variations in connectivity can change both the type of transition and the critical indices. For a polymer substrate modeled by self-avoiding walks (SAWs), the coil--globule transition competes with or synchronizes with magnetic ordering, producing a nontrivial phase diagram \cite{Garel_1999,PhysRevE.104.024122,PhysRevE.108.L042502}. For polymer systems on SAWs without spins, the core geometric and tricritical properties are well studied \cite{Rensburg,Vanderzande_1998,Scaling_concepts}.

Most existing work on spin models with dynamic geometry focuses on discrete spins and has already identified both separated and joint structural--magnetic transitions \cite{Garel_1999,PhysRevE.104.024122,faizullina2021critical1}. Additional scenarios have been demonstrated on deterministic/fractal substrates and on trees (Bethe lattices), with multi-phase diagrams \cite{PhysRevE.106.024130,PhysRevE.108.L042502}, as well as in limiting cases of sparse and compact polymers \cite{papalerosa}. Biophysical models of chromatin based on magnetic polymers (including BEG-type models) further confirm the richness of phase behavior and the presence of multicritical points \cite{PhysRevX.6.041047,PhysRevE.100.052410,D5SM00212E}.

For continuous degrees of freedom, such as the XY model on SAWs, systematic phase-diagram studies have not been carried out. This raises the question of how continuous symmetry affects the location and order of the transitions.


\paragraph*{Aim of the study.}

To investigate the phase behavior of spin systems with dynamic geometry using numerical simulation methods.


\paragraph*{Objectives.}

\begin{enumerate}
	\item To develop a computational framework for spin systems on a regular lattice with dynamically evolving geometry generated by self-avoiding walks (SAWs), with joint sampling of geometric and spin degrees of freedom.
	\item To simulate spin systems with dynamic geometry using the developed framework, and to apply it to the Ising and XY models on SAWs.
	\item To analyze the numerical data and, whenever a phase transition is present, to determine its characteristics (location, order, and critical behavior).
\end{enumerate}


\paragraph*{Original contribution.}

\begin{enumerate}
	\item A computational method is developed and implemented for spin systems on a regular lattice with dynamic geometry defined by self-avoiding walks (SAWs). The implemented approach allows joint averaging over geometric and spin degrees of freedom.
	\item Using the developed method, we simulate systems with discrete (Ising) and continuous (XY) spins.
	\item A data-analysis protocol is formulated for spin systems with dynamically evolving geometry. We characterize behavior on SAWs near phase-transition points for the Ising model.
	\item The XY model on SAWs is investigated for the first time.
\end{enumerate}


\paragraph*{Methodology and methods.}

We employ Markov chain Monte Carlo methods and finite-size analysis for spin models on SAWs. The work uses the developed simulation framework for spin systems on a regular lattice with dynamically evolving geometry generated by SAWs, using joint sampling over polymer conformations and spin configurations and the accompanying analysis protocol tailored to dynamically evolving geometry.


\paragraph*{Theoretical and practical significance.}



The study develops the theoretical understanding of spin systems on dynamically evolving lattices generated by random self-avoiding walks. The results extend current knowledge of the relationship between magnetic and structural phase transitions in polymer systems and make it possible to trace the extent to which critical behavior is inherited from parent models on regular lattices. The practical significance of the work is associated with the development and implementation of numerical methodology for such systems. A computer program enabling computational experiments and data processing was also created within the study.



\section*{Synopsis by chapters}

\noindent\textbf{Introduction.}
We formulate the problem and motivation: to study phase transitions in spin models on dynamic geometry generated by self-avoiding walks on a regular lattice. We provide a literature review to assess the state of the art.

\medskip
\noindent\textbf{Chapter 1. Spin models and dynamic geometry.}
The chapter defines the homopolymer model, represented by self-avoiding walks in two and three dimensions, introduces geometric observables such as the mean radius and radius of gyration and their scaling properties, and discusses the dynamic HP model as a simple polymer model with internal monomer states. The Ising and XY models on dynamic geometry, represented by SAWs, are described. The chapter also formulates the context of phase transitions in polymer and spin systems and gives the expected critical exponents in limiting regimes.

\medskip
\noindent\textbf{Chapter 2. Numerical methods and data analysis.}
The chapter presents the numerical methods used in the dissertation. Markov chain Monte Carlo algorithms for the canonical and grand-canonical ensembles are described for polymer and magnetic-polymer models on SAWs. For selected simulation series, extended-sampling methods are used, including replica exchange (parallel tempering) and multihistogram reweighting. To locate critical points and analyze the character of the transitions, the study employs scaling relations for geometric radii, Binder cumulants, pairwise regressions for curve-intersection points, and analysis of the shape of energy distributions.

\medskip
\noindent\textbf{Chapter 3. Dynamic HP model.}
This chapter investigates the dynamic HP model on the square lattice, in which both the conformations of the polymer chain and the internal states of the monomers are dynamical. Based on Markov chain Monte Carlo simulations in the grand-canonical ensemble, the phase diagram of the model is constructed and its geometric properties are analyzed as the interaction parameter is varied. The study is qualitative in nature and serves as a basis for the subsequent analysis of more complex models of magnetic polymers.

\medskip
\noindent\textbf{Chapter 4. Ising model on SAWs (2D).}
We present simulation results and their analysis for the Ising model on SAWs. The coil--globule transition is continuous and inherits the universality class of the $\Theta$-polymer; numerically, $J_{\Theta}\approx 0.833(1)$. The magnetic transition is also continuous; $J_{c}\approx 0.8340(5)$. Energy distributions remain unimodal near both transitions. The near-coincidence of $J_{\Theta}$ and $J_{c}$ within errors indicates almost co-located geometric and magnetic criteria while maintaining a continuous character for both.

\medskip
\noindent\textbf{Chapter 5. XY model on SAWs.}
We present simulation results and analysis for the XY model on SAWs. 
\emph{Two dimensions:} the coil--globule transition is continuous and consistent with the $\Theta$-polymer universality class; $J_{\Theta}\approx 1.3634\pm 0.005$. The magnetic transition occurs at larger $J$ and is continuous; $J_{c}\approx 1.435\pm 0.008$. Energy distributions show no bimodality near either transition, consistent with the absence of latent heat.\\
\emph{Three dimensions:} the structural transition is localized near $J_{\Theta}\approx 0.876(5)$. The scaling of the radii remains compatible with the polymer value $\nu_{\Theta}=1/2$; however, the distributions of geometric observables and the behavior of the Binder cumulant for the number of topological contacts show that the structural rearrangement cannot be reduced to a simple continuous-transition picture. The magnetic transition exhibits first-order signatures: bimodal distributions of the energy per spin and a negative dip of the Binder cumulant in a narrow interval of $J$.

\medskip
\noindent\textbf{Conclusion.}
The dissertation summarizes the main results. For the dynamic HP model, the phase behavior is studied qualitatively and the region of the geometric transition is localized. For spin models on SAWs, geometric universality is confirmed when Ising and XY spins are added. In two dimensions, the Ising and XY models exhibit continuous transitions; in the XY case, the geometric and magnetic criteria are separated. In the three-dimensional XY model, robust numerical signatures of a first-order magnetic transition are found, together with strong coupling between magnetic ordering and chain compactness. The numerical approach and data-analysis methods provide stable and reproducible results.


\paragraph*{Theses submitted for defense.}

\begin{enumerate}
\item A Markov chain Monte Carlo algorithm has been developed for the numerical simulation of spin systems with dynamically evolving geometry defined by self-avoiding walks. This approach makes it possible to study spin models on random polymer trajectories in which geometric and magnetic degrees of freedom evolve jointly.
\item The Ising model on dynamic structures generated by SAWs on a regular lattice has been studied numerically. The data show that, in two dimensions, the structural transition between globular and coil states of the polymer chain and the magnetic transition occur in the same parameter region within numerical uncertainty. On the square lattice, the magnetic transition is continuous. Estimates of the metric critical exponent characterizing the scaling of the polymer-chain size in the coil--globule transition region indicate inheritance from the parent model, the classical homopolymer.
\item The XY model on self-avoiding walks on square and cubic lattices has been studied numerically for the first time. A phase transition between globular and coil states of the polymer chain is observed. In two dimensions, the geometric transition is consistent with the parent homopolymer model, whereas in three dimensions additional indicators point to a sharp structural rearrangement. The magnetic phase transition is continuous in two dimensions and demonstrates first-order signatures in three dimensions.
\end{enumerate}


\paragraph*{Personal contribution of the author.}

The problem formulation was proposed by E.~A.~Burovskiy. Hypotheses, ideas, and the method were developed jointly. The implementation and optimization of the numerical simulations, as well as the data analysis, were performed by the author. In the publications, the author’s contribution is decisive. Conference presentations were delivered by the author.

During the preparation of the dissertation text, L.~N.~Shchur provided scientific guidance: advising on the formulation of goals, selection and validation of analysis methods, structure of presentation, and interpretation of results.


\paragraph*{General conclusions.}

\medskip
\noindent\textbf{Main results.}
\begin{enumerate}
\item A numerical framework is developed and verified for modeling polymer and spin-polymer systems on self-avoiding walks. The approach is applied to the dynamic HP model and to spin models with discrete (Ising) and continuous (XY) degrees of freedom.
\item For the dynamic HP model on the square lattice in the grand-canonical ensemble, the phase diagram is constructed and the geometric behavior of the chain is studied qualitatively as the interaction parameter is varied. It is shown that, as the interaction is strengthened, the system changes from a regime close to a non-interacting self-avoiding walk to more compact configurations, while the region of the coil--globule structural transition is shifted relative to the homopolymer case.
	\item For the 2D Ising model on SAWs, both the structural (coil--globule) and magnetic transitions are continuous. Their locations coincide within errors; energy distributions are unimodal, consistent with the absence of latent heat.
\item In the two-dimensional Ising and XY models, the geometric transition retains a continuous character and is consistent with the parent polymer model. In the three-dimensional XY model, the metric scaling of the radii remains compatible with $\nu_{\Theta}=1/2$, but additional indicators, namely the Binder cumulant for contacts and distributions of observables, point to first-order signatures. Magnetic behavior depends substantially on dimensionality: in two dimensions the magnetic transitions remain continuous, whereas in the three-dimensional XY model a transition with first-order signatures occurs.
 \item The numerical analysis methods used in the dissertation, including finite-size analysis, Binder cumulants, pairwise regressions of intersection points, and analysis of energy distributions, provide consistent localization of critical regions and reproducible conclusions. For the Ising and XY models, additional data obtained using replica exchange and multihistogram reweighting served as an independent check of the main results.
\end{enumerate}


\subsection*{Dissemination of results.}

\paragraph*{Publications.}

\begin{enumerate}
	\item Faizullina K., Burovski E. \emph{Globule--coil transition in the dynamic HP model} //
	Journal of Physics: Conference Series, 1740 (2021), 012014.
	\item Faizullina K., Pchelintsev I., Burovski E. \emph{Critical and geometric properties
		of magnetic polymers across the globule--coil transition} // Physical Review E 104(5) (2021), 054501.
	\item Faizullina K., Burovski E. \emph{XY model on self-avoiding walks} // Physical Review E 112(2) (2025), 024107.
	\item Faizullina K., Burovski E. \emph{XY model on self-avoiding walks: a large-scale Monte Carlo study} // Supercomputing in Russia (Proceedings), 2022, pp.~24--30.
\end{enumerate}


\paragraph*{Presentations.}

\begin{enumerate}
	\item International workshop ``StatPoly: Polymers and Biopolymers'', Venice, Italy, 20--23 July 2025, poster: \emph{Long-Range XY Model on Self-Avoiding Walks}.
	\item XXXV IUPAP Conference on Computational Physics (CCP 2024), Thessaloniki, Greece, 7--12 July 2024, oral: \emph{XY model on a self-avoiding walk}.
	\item ``Supercomputing in Russia'', Moscow, Russia, 26--27 September 2022, oral: \emph{XY model on self-avoiding walks: a large-scale Monte Carlo study}.
	\item 33rd IUPAP Conference on Computational Physics (CCP 2022), Online, 1--4 August 2022, oral: \emph{XY model on self-avoiding walks}.
\end{enumerate}


\paragraph*{Registered software.}

\begin{enumerate}
	\item \emph{Search for a minimum-energy conformation in the hydrophobic--polar model} (Certificate of state registration of a computer program No.~2019667328, dated 23.12.2019).
\end{enumerate}
